%%%%%%%%%%%%%%%
%
% $Autor: Wings $
% $Datum: 2020-01-29 07:55:27Z $
% $Pfad: General/IMU.tex
% $Version: 1785 $
%
%
%%%%%%%%%%%%%%%


\chapter{Inertial Measurement Unit LSM9DS1}

\textcolor{red}{Quellen für die Bilder fehlen!}


\section{Allgemeine Informationen zu einer inertialen Messeinheit}

IMU steht für \textbf{I}nertial \textbf{M}easurement \textbf{U}nit, also eine inertiale Messeinheit. Eine IMU ist eine Kombination aus zwei oder drei MEMS. Zu den Einzelsensoren einer IMU gehören der Beschleunigungs- und Drehratensensor (Gyroskop) und ein Magnetometer. Jeder der Sensoren misst in den drei Richtungen $X$, $Y$ und $Z$. Mit der Anzahl der Sensoren und den Messrichtungen wird der Freiheitsgrad angegeben. Er bestimmt sich aus der Anzahl der Sensoren multipliziert mit den Richtungen: 

\begin{center}
  3 Sensoren $\times$ 3 Richtungen = 9 Freiheitsgrade. 
\end{center}


Die Freiheitsgrade werden im Englischen als \glqq Degrees of Freedom\grqq{} (DOF) bezeichnet. Oftmals sind alle drei Sensoren zusammen in einem Mikrochip untergebracht, dies ist aber nicht zwangsläufig der Fall. Im Mikrochip selbst befindet sich außerdem noch ein Mikrocontroller, damit die Rohdaten über eine Schnittstelle ausgelesen und verarbeitet werden können. Häufig unterstützt eine IMU mehrere Schnittstellen. Das Anbinden erfolgt dann über Inter-Integrated Circuit (I2C), Serial Peripheral Interface (SPI) oder Universal Asynchronous Receiver Transmitter (UART). \cite{Weber:2021}



\subsection{Beschleunigungssensor}

Beschleunigungssensoren messen eine durch Beschleunigung ausgeübte Kraft auf eine Masse. Mit ihnen lassen sich die Größen Kraft, Geschwindigkeit und Weg bestimmen. Die am häufigsten verwendeten Beschleunigungssensoren sind die kapazitiven Beschleunigungssensoren. \cite{Traenkler:2014} Bei diesen Sensoren befindet sich an zwei Federn eine gelagerte Masse, welche nach rechts und links frei beweglich ist. Erfährt der Sensor eine Beschleunigung, verändert sich der Abstand der Kondensatorplatten (C1 und C2) durch die Trägheit der Masse. 

\begin{figure}[H]
	\centering
	\includegraphics[width=\textwidth]{IMU/Kapazitaet}
	\caption{Funktionsprinzip eines kapazitiven 1D Beschleunigungssensors}
\end{figure}



Stärke und Richtung der Auslenkung werden anschließend mittels einer kapazitiven Messung bestimmt. Für jede Achse wird eine solche Messung durchgeführt, insgesamt also dreimal. Auf einem speziellen Chip werden anschließend die Daten der Messung aufbereitet. \cite{Weber:2021}


\subsection{Drehratensensor}

Ein Drehratensensor (Gyroskop) misst die Rotationsgeschwindigkeit um eine definierte Achse.  Weit verbreitet ist das Messprinzip des Zweimassenschwingers, welches auf der Corioliskraft beruht. Hier schwingen zwei seismische Massen, die über eine Feder verbunden sind, gegensinnig radial, solange keine Rotation vorliegt. Erfährt der Sensor eine Beschleunigung, bewirkt die Corioliskraft eine Auslenkung der beiden Massen. \cite{Traenkler:2014} 
Wie zuvor bei dem Beschleunigungssensor bewirkt eine Auslenkung eine Veränderung des Verhältnisses der Kondensatorplatten.  Die Kapazität an der einen Platte erhöht sich, während sich die Kapazität an der anderen verringert.  Mittels kapazitiver Messung werden Richtung und Stärke der Auslenkung bestimmt. Wie bei dem Beschleunigungssensor werden die Messungen für jede Achsenrichtung durchgeführt. 


\subsection{Magnetometer}

Magnetometer messen das Magnetfeld für jede Achse. Hierfür wird in einer IMU häufig der Hall-Effekt verwendet. Dieser Effekt beschreibt, dass sich an einem stromdurchflossenen Leiter, welcher sich in einem Magnetfeld befindet, eine Spannung messen lässt. Diese messbare Spannung wird Hallspannung genannt. Die Stärke der Hallspannung ist abhängig von dem fließenden Strom, dem Magnetfeld und von dem Leitermaterial.

\begin{figure}[H]
	\centering
	\includegraphics[width=\textwidth]{IMU/Hallsensor}
	\caption{Hallspannung}
\end{figure}



Bei kapazitiven Magnetometern wird eine drehbare Masse mittels Federn gelagert. Ein Magnetfeld bewirkt auch hier eine Auslenkung und damit eine messbare Kapazitätsänderung.

\begin{figure}[H]
	\centering
	\includegraphics[width=0.7\textwidth]{IMU/Magnetometer}
	\caption{Funktionsprinzip eines kapazitiven Magnetometer}
\end{figure}



Dies ist bereits am Magnetfeld der Erde messbar. Wenn der Sensor beispielsweise zum Nordpol ausgerichtet wird, dann zeigt die Achse, welche dem Pol am nächsten ist, ein größeres Magnetfeld an, als die anderen beiden. Magnetometer sind demnach also sehr empfindlich und daher sehr anfällig für Störungen. Bei dem Einbau von Magnetometern muss deshalb darauf geachtet werden, dass man sie nicht in der Nähe von Motoren oder anderen Geräten, die ein magnetisches Feld erzeugen, verbaut. Andernfalls geben sie aufgrund der Störungen fehlerhafte Werte aus. Alle genannten Sensoren unterliegen Drift, auf den in Kapitel 5.5 näher eingegangen wird. \cite{Weber:2021}




\section{Spezielle Betrachtung der verwendeten IMU}

Bei der verwendeten IMU handelt es sich um das Modell LSM9DS1 von der Firma STMicroelectronics International NV. Sie gehört der Produktfamilie iNEMO (integrated Nano Electro Mechanical Systems) an. Der LSM9DS1 findet unter anderem Einsatz in der Navigation, der Bewegungs- und Gestensteuerung sowie in der Unterhaltungselektronik. Auf dem Chip befinden sich die Schnittstellen I2C und SPI, sowie ein 3-Achsen-Beschleunigungssensor, ein 3-Achsen-Drehratensensor und ein 3-Achsen-Magnetometer.  Die IMU hat also 9 Freiheitsgrade. 

Allerdings konnten dem Datenblatt die verwendeten Messprinzipien der einzelnen Sensoren nicht entnommen werden. \cite{STMicroelectronics:2015}




\section{Kalibrierung der IMU}

Jeder Sensor der IMU wird einzeln kalibriert. Da Mikrosensoren sehr empfindlich auf Vibrationen reagieren, ist es wichtig, diese zu vermeiden. 


\subsection{Beschleunigungssensor}

Um den Beschleunigungssensor zu kalibrieren wird er wie ein Würfel auf seine sechs Seiten gedreht. Durch diese Vorgehensweise werden die Offset-Werte bestimmt, die anschließend mit den gemessenen Werten verrechnet werden. Auf diese Weise kann außerdem ein schiefer Einbau korrigiert werden.


\subsection{Drehratensensor}

Die Kalibrierung des Drehratensensors ist die einfachste. Für einen Nullabgleich der Achsen muss der Sensor nur in den Ruhezustand versetzt werden. Falls das nicht der Fall sein sollte, müssen auch hier die Offset-Werte bestimmt werden.


\subsection{Magnetometer}

Magnetometer werden kalibriert, indem man sie in alle möglichen Richtungen dreht. Häufig wird bei der Variante der Sensor in einer Schleifenbewegung bewegt, beziehungsweise wird versucht, mit dem Sensor eine ‚Acht‘ zu zeichnen. Alternativ kann man aber auch wie beim Beschleunigungssensor das Magnetometer wie ein Würfel auf jede Seite kippen. \cite{Weber:2021}

\clearpage


\section{Kalibrierung der speziellen IMU}

Das Problem bei dem Arduino Winkelmesser ist, dass er im Ruhezustand auf einer ebenen Oberfläche stark driftet. Das Driften äußert sich darin, dass Winkel angezeigt werden, obwohl der Winkelmesser auf einer ebenen Oberfläche steht. Sobald man die IMU allerdings neigt, werden die ausgegebenen Werte der Winkel realistisch. Der Hersteller STMicroelectronics International NV hat zwar eine Anleitung wie man den LSM9DS1 kalibriert bereitgestellt, aber das betrifft jedoch nur die IMU ohne den Arduino. \cite{STMicroelectronics:2022}


In den Bibliotheken, die wir für das Programm nutzen, gibt es keine Kalibrierfunktion. Nachdem erhebliche Probleme in der Kalibrierung der IMU aufgetreten sind, hat die Gruppe Kontakt mit dem Arduino-Kundenservice aufgenommen. Dem war dieses Problem bereits bekannt, und klärte uns darüber auf, dass die Bibliotheken über keine Kalibrierfunktion verfügen. \vspace{0.5cm}

\begin{figure}[H]
	\centering
	\includegraphics[width=0.45\textwidth]{IMU/SupportEmail}
	\caption{Antwort des Arduino Support Teams}
\end{figure}

Die einzige Möglichkeit einer Kalibrierung besteht darin, dass man diese im Programm durchführt. Abhilfe für unser Problem schafften wir also im Sketch, indem für alle Winkel die kleiner gleich 10 Grad sind, 0 Grad auf der Anzeige ausgegeben wird.




\section{Drift und Ungenauigkeiten}

Die Aufgabe eines Sensors ist es, eine physikalische, chemische oder biologische Messgröße in elektrische Signale umzuwandeln, damit diese ausgewertet werden können. Es kommt allerdings vor, dass sich die Messwerte ändern, obwohl die gemessene Größe sich selbst nicht verändert. Dieses Phänomen der Abweichung wird als Drift bezeichnet. \cite{Traenkler:2014} 

Das Driften von Sensoren ist demnach ein unerwünschter Effekt. Er ist zeitabhängig und tritt im Sensor bedingt durch Alterung und durch Ungenauigkeiten auf. Im schlimmsten Fall können diese Veränderung den Sensor so sehr beeinträchtigen, dass es zu einem funktionalen Ausfall kommt.\cite{Kern:2023} 

Außerdem spielen Umgebungsbedingungen eine ebenso wichtige Rolle, da diese die Messabweichungen zusätzlich verstärken.
Ein häufiges Beispiel ist hierfür das thermische Driften, welches durch Temperaturschwankungen verursacht wird. \cite{Georgy:2010}
Neben den Temperaturschwankungen beeinflussen weitere Fehlerursachen wie eine Veränderung der Empfindlichkeit, des Nullpunktes, der Frequenz sowie 
logische Fehlfunktionen und Rauschen das Driften. \cite{Traenkler:2014}
Bis auf das Rauschen, welches den zufälligen Messfehlern zugeordnet wird, lassen sich die anderen kompensieren und gehören demnach den systematischen Messfehlern an. \cite{Leon:2019}

\begin{figure}[H]
	\centering
	\includegraphics[width=0.8\textwidth]{IMU/drift.png}
	\caption{Auswirkungen von Drift auf Messergebnisse}
\end{figure}

 Um die Einflüsse der systematischen Messfehler und damit Drift zu verringern, werden Sensoren regelmäßig nachkalibriert. 
Dennoch kann es geschehen, dass ein Sensor, welcher zwischen zwei Kalibrierungen betrieben wird, driftet. \cite{Traenkler:2014}
Ein weiterer Grund für Abweichung sind die zufälligen Messfehler oder auch stochastische Messfehler.
Stochastische Fehler treten beispielsweise durch Rauschen auf und führen, falls sie nicht korrigiert werden, ebenfalls zu Drift.
Ein übliches Verfahren diese Fehler zu minimieren ist der Kalman Filter. \cite{Narasimhappa:2020}
Der Nachteil des Kalman Filters ist allerdings, dass er während des Linearisierungsprozesses unter Divergenz aufgrund von Ausfällen leidet.
Aufgrund ihrer charakteristischen Fehler und der stochastischen Natur sind MEMS-Inertialsensoren mittels Kalman Filter nur begrenzt modellierbar.
Im Gegensatz dazu hat sich der Partikel-Filter als sehr nützlich erwiesen, da er die Fähigkeit besitzt, 
auch mit nichtlinearen und nicht-Gauß'schen Modellen umzugehen. \cite{Leon:2019}

Eine weitere in der Praxis gängige Möglichkeit die Ungenauigkeiten von Sensoren zu verringern ist die sogenannte Sensorfusion. Hier wird ein zusätzlicher Sensor verwendet, um die Probleme eines anderen Sensors zu kompensieren.  Ein Beispiel für eine solche Verwendung findet sich in einer Drohne. Da das Magnetometer der IMU Drift unterliegt, verändert sich die Genauigkeit mit der Zeit. Diese Genauigkeitsänderung kann zur Folge haben, dass die Drohne vom Kurs abkommt. Um eine Kursänderung zu verhindern, wird ein GPS-Sensor verwendet. Das GPS kann im Stillstand allerdings keinen Kurs berechnen. Umgekehrt kann also die IMU verwendet werden und so den Nachteil des GPS-Sensors ausgleichen. \cite{Weber:2021}



\section{Algorithmen zur Auswertung einer IMU}

Kalman-Filter (KF) werden verwendet, um in Echtzeit die Daten dynamischer Systeme auf niedrigem Signalbearbeitungsniveau zu fusionieren. In vielen Anwendungen, wie zum Beispiel bei Navigations- und Führungssystemen, bei der Radarverfolgung und bei Abstandmessungen hat sich der Kalman-Filter als sehr nützlich erwiesen. Um die statistisch wahrscheinlichsten Sensorwerte zu bestimmen, wird ausgehend von einem linearen Systemmodell und einer Gauß'schen Fehlerverteilung für System und Sensor, ein rekursives Schätzverfahren verwendet.  Dieses Verfahren zeichnet sich besonders durch die hohe Verarbeitungsgeschwindigkeit und die geringe Speicheranforderungen aus. \cite{Traenkler:2014} 

Wie schon im Kapitel ‚Drift‘ angemerkt, sind Inertialsensoren nicht linear, sondern von stochastischer Natur und deshalb mit dem Kalman-Filter schwer zu modellieren. Bei seiner Verwendung werden linearisierte Modelle der Zustände verwendet und dabei nur Terme erster Ordnung der Taylorreihenentwicklung durch den Kalman-Filter berücksichtigt.  Dies führt dazu, dass der Filter die Navigationszustände und nicht die Zustände selbst schätzt. Die Auswertung einer IMU mittels Kalman-Filter ist aus diesem Grund nur bedingt geeignet. \cite{Georgy:2010} 


Für eine solche Anwendung ist der erweiterte Kalman-Filter (EKF) der passendere Algorithmus. Bei diesem Auswertealgorithmus handelt es sich ebenfalls um ein Schätzverfahren, welches Näherungsmodelle für die Zustände benutzt. Hierbei muss beachtet werden, dass die Anfangszustände nahe bei den wirklichen Zuständen liegen. Ist das nicht der Fall, wird die Schätzung bedeutungslos. Um sicherzustellen, dass die Terme zweiter und höherer Ordnung der Taylorreihenentwicklung unbedeutend sind, müssen während der Filterung die Vorhersagen nah bei den wahren Zuständen liegen. Sollte der gewählte Nominalzustand zu weit vom wahren Wert entfernt liegen, fällt die tatsächliche Kovarianz größer aus als die Geschätzte und der Filter passt nicht mehr zum System. Der EKF ist im Vergleich zum Kalman-Filter der leistungsfähigere Algorithmus, von Nachteil ist allerdings der erhöhte Rechenaufwand und die Komplexität des erweiterten Kalman-Filters. Ein weiterer Algorithmus um mit nichtlinearen und nicht-Gauß‘schen Modellen umzugehen ist der Partikel-Filter. Der Partikel-Filter ist außerdem in der Lage, die Sensorcharakteristiken, die Bewegungsdynamik und die Rauschverteilung anzupassen. \cite{Georgy:2010} 


Nachteil des Partikelfilters ist jedoch, dass die benötigte Partikelzahl systematisch schwer zu wählen ist und sich durch die steigende Partikelzahl der Rechenaufwand erhöht. \cite{Traenkler:2014}


\section{Algorithmen zur Auswertung der speziellen IMU}

Wie schon zuvor bei der Kalibrierung fanden wir auch für die Auswertungsalgorithmen des LSM9DS1 keine weiteren Informationen. Daran änderte auch eine Nachfrage bei Arduino nichts. Da in der verwendeten IMU nichtlineare Sensoren verbaut sind, ist es wahrscheinlich, dass für die Auswertung der Erweiterte Kalman-Filter verwendet wird. Verwendung, sowie vor und Nachteile des EKF wurden im vorherigen Kapitel bereits erläutert.

